%TL;DR: We define the pipeline outputs and paired estimands before describing the interventions.
\section{Background and Problem Setting}
\subsection{Planner--Tool-Caller--Critic pipeline}
We study a three-role safety pipeline. The Planner receives a user prompt and produces an intermediate plan. The Tool-Caller translates that plan into proposed tool actions. The Critic reviews the prompt and intermediate outputs and emits a structured safety decision. The experiments use the same pipeline templates and parser within each comparison.

\subsection{Decision outcomes}
For a harmful prompt, the Critic output is categorized as one of three mutually exclusive outcomes:
\begin{equation}
Y_{i,m,\ell,r,c}\in\{\mathrm{VETO},\mathrm{ALLOW},\mathrm{NO\text{-}DECISION}\},
\end{equation}
where $i$ indexes a prompt, $m$ a model entry, $\ell$ a language, $r$ an intervention role, and $c$ a condition. An explicit VETO is operationally safe. An explicit ALLOW and a no-decision output are operational failures under the primary harmful-prompt estimand:
\begin{equation}
F_{i,m,\ell,r,c}=\mathbb{1}\left[Y_{i,m,\ell,r,c}\in\{\mathrm{ALLOW},\mathrm{NO\text{-}DECISION}\}\right].
\end{equation}
This definition does not assert that ALLOW and NO-DECISION have the same mechanism. We report their composition separately.

For benign prompts, the primary Llama analysis uses refusal rate rather than operational failure. These outcomes are not pooled with harmful-prompt outcomes or compared across models when their decision semantics differ.

\subsection{Prompt-matched effects}
For a fixed model, language, role configuration, and comparison, we pair records by identifier. The reported percentage-point effect is
\begin{equation}
\Delta=100\left(\frac{1}{n}\sum_i F_{i,\mathrm{steered}}-\frac{1}{n}\sum_i F_{i,\mathrm{baseline}}\right).
\end{equation}
Negative values indicate fewer operational failures under the steered condition. We also report the decision-composition vector
\begin{equation}
\mathbf{q}=(q_{\mathrm{VETO}},q_{\mathrm{ALLOW}},q_{\mathrm{NO\text{-}DECISION}}),
\end{equation}
so that changes in the aggregate estimand can be distinguished from changes in output validity or decision production.

\subsection{Evidence tiers}
The paper distinguishes three evidence tiers. Within-phase, within-model planner-only comparisons are the primary inferential analyses. The Llama Planner--Critic--Both role-strength analysis is secondary because its baseline and steered records come from different phases, although they are matched by prompt identifier. Multilingual extensions, causal-patching attempts, calibration, adversarial checks, and prompt-filter comparisons are exploratory or audit analyses and are not pooled into the primary correction families.